\subsubsection{Architecture}
The ForGAN architecture described in Sec.~\ref{sec:forgan} was used, without modifications. In $G$, the condition passes through an LSTM layer. The final hidden state is concatenated with the noise vector $z$. The result passes through a dense layer of width $R_G+N$ , then a ReLU activation and another dense layer of width one, where $R_G$ is the  LSTM hidden size and $N$ is noise dimension.

In $D$, the candidate value ($\hat{x}_{t+1}$ or $x_{t+1}$) is appended to the condition, that gives a sequence of size $l+1$. The result passes through an LSTM layer with the hidden size $R_D$. The final hidden state passes through a dense layer with a sigmoid activation.

Following the the implementation of \textcite{vuleticFinGANForecastingClassifying2024} $N=8$, $R_G=8$ and $R_D$=8. However $N=8$, $R_D=8$ are varied in Sec.~\ref{sec:configurations}.



\subsubsection{Training configuration}

\subsubsection{Adaptations}

\todo{The implementation is based on the reference Fin-GAN code published by \textcite{vuleticFinGANForecastingClassifying2024}. The network definitions, the gradient norm matching procedure and the training loops were adopted from that implementation.}